\documentclass[10pt,a4paper]{amsart}
\usepackage[margin=19mm]{geometry}
\usepackage{amsmath,amssymb,mathtools}
\usepackage[scr=rsfs,cal=euler]{mathalfa}
\usepackage{xfrac}
\usepackage[unicode,hidelinks,bookmarks=false]{hyperref}
\newcommand{\ie}{i.e.\ }
\newcommand{\cf}{cf.\ }
\newcommand{\resp}{respectively\ }
\newcommand{\Subsection}{\subsection}
\providecommand{\nobreakdash}{\nobreak\hbox{-}}
\setlength{\emergencystretch}{2em}
\multlinegap0pt
\usepackage{amssymb}
\usepackage{mathtools}
\usepackage{tikz}
\newtheorem{lem}{Lemma}
\newtheorem{cor}{Corollary}
\usepackage{cleveref}
\crefname{lem}{Lemma}{Lemmas}
\crefname{cor}{Corollary}{Corollarys}
\DeclareMathOperator{\sink}{sink}
\title{Poset topology, moves, and Bruhat interval polytope lattices}
\author{Christian Gaetz, Patricia Hersh}
\date{Source: arXiv:2410.08076v2 (CC BY 4.0)}
\begin{document}
\maketitle

\noindent\emph{Source and licence.} Poset topology, moves, and Bruhat interval polytope lattices. Version arXiv:2410.08076v2 (CC BY 4.0), distributed by arXiv under the Creative Commons Attribution 4.0 International licence, http://creativecommons.org/licenses/by/4.0/, as recorded in the arXiv licence metadata for this exact version and shown on the abstract page https://arxiv.org/abs/2410.08076v2. Full licence notice: \texttt{licenses/arxiv-2410.08076-CC-BY-4.0.txt}.\par\smallskip

\noindent\textit{Editorial scope.} Counted unit: the article's lem lem:vertex-disjoint (source0.tex lines 653--655). The article's complete original proof of it is delivered with it (source0.tex lines 657--664, one \texttt{proof} environment of 8 lines). No mathematics is written, repaired or re-derived: the statement and the proof are the article's own lines, in the article's own notation, and no retained line is substituted or re-flowed.  Retained uncounted as the source-local context the proof needs: lines 516--518; lines 533--535; lines 544--551; lines 641--645.  Nothing was omitted from any retained span, so the package carries no omission record.  No retained line contains a citation, so the wrapper carries no bibliography and no bibliography entry.  No retained line contains a figure environment, an image-inclusion command or drawing code, so no image is delivered and the package declares no asset dependency.  Editorial formatting only, in addition: an AMS \texttt{amsart} preamble that restates the article's own declarations and macros used by the retained lines, namely the environments \texttt{lem}, \texttt{cor} on separate counters, together with the standard packages those lines need.  The retained environments print in this document as Lemma 1, Corollary 1 in order of appearance, whereas the article prints the same items with the numbering of its own full document.  The complete native source of the article as distributed in the arXiv e-print for this exact version is bundled unchanged as \texttt{sources/166635/source0.tex}.
\begin{lem}\label{lem:path-exists}
Let $Q$ be an $\mathcal{O}$-directionally simple polytope such that $\mathcal{O}$ is the Hasse diagram of a lattice $L$. Then, for any $u \leq_L v$, and any chains $\gamma_1,\gamma_2 \in \mathcal{M}(u,v,\mathcal{O})$, there exists a path from $\gamma_1$ to $\gamma_2$ such that all intermediate chains $\gamma$ contain all vertices in $\gamma_1 \cap \gamma_2$. 
\end{lem}

\begin{cor}\label{cor:sink-lemma}
In the notation of \Cref{lem:path-exists} and its proof, if $\sink(F)\in \gamma_1$, then all elements of $\gamma_1 * \gamma_2$ other than $\gamma_1$ contain the edge $e_2$.
\end{cor}

\begin{lem}\label{lem:unique-sequence}
Let $Q$ be an $\mathcal{O}$-directionally simple polytope such that $\mathcal{O}$ is the Hasse diagram of a lattice $L$.   Suppose  $\gamma $ and $\gamma_1$ are saturated chains in an interval $[u,v]_L$ that include distinct cover relations upward from $u$.  Let $\gamma_2$ be any node, distinct from $\gamma $, 
on the path $\gamma * \gamma_1$.   Denote by $e_1$  the lowest edge in $\gamma_1$, and denote by  $e_2$ the lowest edge in $\gamma_2$ that is  not in $\gamma$.  
  
Then there exists an  alternating sequence $(\epsilon_0,F_1,\epsilon_1,F_2,\dots ,F_r,\epsilon_r)$ such that  $\epsilon_0 = e_1$, $\epsilon_r = e_2$, the edges $\epsilon_0,\dots ,\epsilon_r$ are distinct, 
each  $\epsilon_i$ is an edge of $Q$ whose lower vertex $x_i$ is in $\gamma$,  and
 each $F_i$ is a 2-face of $Q$ containing  $\epsilon_{i-1}, \epsilon_i$ and  all edges and vertices of $\gamma$ between $x_{i-1}$ and $x_i$.
\end{lem}

\begin{lem}\label{lem:what-determines}
 The sequence $(\epsilon_0,F_1,\epsilon_1,F_2,\dots ,F_r,\epsilon_r)$ from  
 \Cref{lem:unique-sequence}
  is uniquely determined by $\gamma $ and $e_2$. In particular, $\gamma$ and $e_2$ determine $e_1$. 
\end{lem}

\begin{lem}\label{lem:vertex-disjoint}
Let $Q$ be an $\mathcal{O}$-directionally simple polytope such that $\mathcal{O}$ is the Hasse diagram of a lattice $L$.   If $\gamma, \gamma_1$, and $\gamma_1'$ are saturated chains in an interval $[u,v]_L$ using distinct edges upward from $u$, then $(\gamma * \gamma_1) \cap (\gamma * \gamma_1')=\{\gamma\}$.
\end{lem}

\begin{proof}
Suppose that $\gamma \neq \gamma '' \in (\gamma * \gamma_1)\cap (\gamma * \gamma_1')$.
Since $\gamma_1$ and $\gamma_1'$ have distinct lowest edges, $\gamma ''$ must have the same lowest edge as $\gamma$. Letting $F$ (resp. $F'$) denote the unique 2-face containing the lowest edges in $\gamma $ and $\gamma_1 $ (resp. $\gamma $ and $\gamma_1'$), it follows that $\gamma''$ must occur prior to the flip across $F$ (resp. $F'$) within $\gamma * \gamma_1$ (resp. $\gamma * \gamma_1'$). If $\sink (F)\in \gamma $ this would contradict \Cref{cor:sink-lemma}. Thus it suffices to consider the case
$\sink (F) \not \in \gamma $ (resp. $\sink (F') \not \in \gamma $). 
This implies that $r\ge 1$ in  the sequence $(\epsilon_0,F_1,\dots ,\epsilon_r)$ determined according to \Cref{lem:unique-sequence} by $\gamma$ and the lowest edge in $\gamma ''$ not in $\gamma$. This allows us to  assume $r>0$ in the remainder of this proof.

Consider any node $\gamma_2$ on the path $\gamma*\gamma_1$ and any node $\gamma_2'$ on the path $\gamma * \gamma_1'$.  Let $e_2$ be the lowest edge in $\gamma_2$ not shared with $\gamma $, and let $e_2'$ be the lowest edge in $\gamma_2' $ not shared with $\gamma $.   To prove $\gamma_2 \ne \gamma_2'$, it suffices to prove $e_2 \ne e_2'$.  If $e_2 = e_2'$, then \Cref{lem:what-determines} shows that $\gamma $ and $e_2$ would give rise to the same  alternating sequence of edges and 2-faces that  $\gamma $ and $e_2'$ would produce.  In particular, they would give rise to the same edges upward from $u$, contradicting $\gamma_1$ and $\gamma_1'$ having  distinct edges upward from $u$.
\end{proof}

\end{document}
